\scanpage{100}
\mspage{87}

\uline{\deleted{\unsure{Corollaire}}}\inserted{Proposition} 5.2. Soit $(M,W,C)$ une catégorie à cofibrations spéciale. \deleted{Pour que \unsure{le|la}} Alors avec les notations de (5.1), on a
\[
 W=\mathrm{TF}\circ\mathrm{TC} \tag*{(5.8)}
\]
\[
 \mathrm{TF}=F\cap W. \tag*{(5.9)}
\]

\uline{Dém} Soit $f\in W$, écrivons $f=pi$, avec $i\in C$, $p\in\mathrm{TF}$, donc $p\in W$, donc $i\in W$, donc $i\in C\cap W=\mathrm{TC}$, d'où (5.8).

Prouvons (5.9). On a $\mathrm{TF}\subset F\cap W$, prouvons inversement qu'une $f\in F\cap W$ est dans $\mathrm{TF}$.
\diagram{\begin{tikzcd}[ampersand replacement=\&]  X \arrow[rr,"f\in F\cap W"] \arrow[dr,"i\in\mathrm{TC}"'] \& \& Y\\  \& X'\arrow[ur,"p\in\mathrm{TF}"'] \&
\end{tikzcd}}
Écrivons $f=pi$, avec $p\in\mathrm{TF}$, $i\in\mathrm{TC}$. Comme $f\in F$, $\deleted{\unsure{c}}i\in\mathrm{TC}$, et $\mathrm{TC}\longleftrightarrow F$ par définition de $F$, $\exists\,r:X'\unsure{\mathbin{-}|\to}X$ tel que $ri=\mathrm{id}_X$, $fr=p$. Donc $f$ est facteur direct de $p\inserted{\in\mathrm{TF}}$, et comme $\mathrm{TF}=C_*$  est stable par facteurs directs, on conclut que $f\in\mathrm{TF}$.

\uline{Corollaire} 5.3. Soit $(M,W,C)$ une catégorie à cofibrations spéciale. Pour que $(W,C,F\overset{\mathrm{déf}}{=}(\mathrm{TC})_*)$ soit un triple de Quillen dans $M$, il f. et il suffit que \deleted{$\unsure{\mathrm{TC},F|\mathrm{TC},E}$} le couple $(\mathrm{TC},F)$ satisfasse à l'axiome de factorisation, (i.e. que ce soit un précouple de Quillen). Pour que cette structure de Quillen soit \unsure{clos|close}, il f. et il s. que $C$ \deleted{et $W$} \unsure{soit|soient} stable par par facteurs directs.
% archive sheet number: 99
